% Chapter 1: The Configuration Hierarchy
\chapter{The Configuration Hierarchy}
\label{ch:configuration}

Claude Code is not just a chatbot with file access.
It is a configurable development environment with a layered context system
that determines how Claude understands your project, what it can do,
and how it behaves.
This chapter covers the four mechanisms that shape that behavior:
\code{CLAUDE.md}, rules, settings, and output styles.

\section{CLAUDE.md --- Your AI's Operating Manual}
\label{sec:claude-md}

\marginnote[Official]{Run \code{/init} to auto-generate a starter CLAUDE.md from your project structure.\sourceurl{https://code.claude.com/docs/en/best-practices}}

\official{The \code{CLAUDE.md} file is the single most important configuration artifact.}
Claude reads it at session start and uses it to understand your project's conventions,
build commands, architecture decisions, and rules of engagement.

\subsection{Hierarchical Loading}

Claude loads \code{CLAUDE.md} files from multiple locations, merging them in order:

\begin{enumerate}
  \item \code{\textasciitilde/.claude/CLAUDE.md} --- global defaults (applies to all projects)
  \item Project root \code{CLAUDE.md} or \code{.claude/CLAUDE.md} --- project-level instructions
  \item Parent and child directory \code{CLAUDE.md} files --- subdirectory-specific overrides
\end{enumerate}

\marginwarning{If your CLAUDE.md exceeds $\sim$500 lines, Claude will silently deprioritize content near the end. Split into rules files.}

\official{Use the \code{@path/to/file.md} import syntax to reference external files
without duplicating content.}

\subsection{What to Include}

\begin{itemize}
  \item Build and test commands Claude cannot guess (e.g., \code{make pilot}, \code{pytest scripts/tests/})
  \item Code style rules that differ from standard conventions
  \item Architecture decisions that affect how Claude writes code
  \item File organization conventions
  \item Testing requirements and coverage targets
\end{itemize}

\subsection{What to Exclude}

\begin{itemize}
  \item What Claude can infer from the codebase (standard language conventions, obvious patterns)
  \item Long explanations --- keep instructions terse and actionable
  \item Aspirational rules nobody follows --- document what IS, not what you wish
\end{itemize}

\begin{tipbox}[Quality Test]
If Claude keeps violating a rule, one of three things is true:
(1) the CLAUDE.md is too long and the rule is getting lost,
(2) the rule contradicts another rule, or
(3) the rule should be a hook instead of an instruction.
\end{tipbox}

\subsection{Hub-and-Spoke Architecture}

\practitioner{For teams managing many repositories, maintain a central ``hub'' repo
with shared patterns referenced via \code{@\textasciitilde/hub/patterns/}.}

This prevents duplication across repos while keeping each project's \code{CLAUDE.md} concise.

\begin{minted}{markdown}
# Project CLAUDE.md (spoke)
@~/shared-patterns/git.md
@~/shared-patterns/testing.md

## Project-Specific
- Build: `make digital`
- Test: `pytest tests/`
\end{minted}

The hub contains canonical patterns for git commits, testing standards, session workflows,
and other cross-cutting concerns. Each spoke references only the patterns it needs.

\practitioner{Classify repos by context engineering maturity:}

\begin{fullwidth}
\begin{tabular}{@{}llp{8cm}@{}}
\toprule
\textbf{Tier} & \textbf{Components} & \textbf{When to Use} \\
\midrule
T0 & None & Prototypes, learning repos \\
T1 & CLAUDE.md only & Small scripts, documentation projects \\
T2 & CLAUDE.md + settings & Standard development projects \\
T3 & + skills, commands, rules & Multi-module projects with established workflows \\
T4 & + hooks, agents, MCP & Projects requiring automation and quality gates \\
T5 & + agent teams, pipelines & Enterprise-scale projects with CI/CD integration \\
\bottomrule
\end{tabular}
\end{fullwidth}

\margintip{Start at T1. Promote when friction justifies it. Most solo projects stabilize at T2--T3.}

% -------------------------------------------------------------------
\section{Rules: Conditional Context Loading}
\label{sec:rules}

\official{Place \code{.md} files in \code{.claude/rules/} with \code{paths:} frontmatter
to control when they load.}\sourceurl{https://code.claude.com/docs/en/best-practices}

Rules without \code{paths:} frontmatter load on \emph{every} session regardless of
which files you are editing. For monorepos and multi-module projects,
this bloats context with irrelevant instructions.

\begin{minted}{markdown}
---
paths: backend/**/*.py
---
# Backend Python Rules
- Use SQLAlchemy 2.0 async syntax
- All endpoints need OpenAPI docstrings
- Run `pytest backend/tests/` after changes
\end{minted}

\practitioner{Conditional rules are essential for any project with more than three
distinct modules.} Without them, a rule about SQL formatting loads
when you are editing React components, wasting context and occasionally
causing Claude to apply the wrong conventions.

\begin{warnbox}[Common Mistake]
Forgetting \code{paths:} frontmatter. The rule loads globally,
Claude's context fills with irrelevant instructions, and important
rules from CLAUDE.md get deprioritized.
\end{warnbox}

% -------------------------------------------------------------------
\section{Settings: 4-Level Precedence}
\label{sec:settings}

\marginnote[Official]{Settings documentation:\sourceurl{https://code.claude.com/docs/en/settings}}

Claude Code uses a 4-level settings hierarchy. Higher levels override lower:

\begin{enumerate}
  \item \textbf{Managed} (\code{.managed-settings.json}) --- IT-deployed, highest priority
  \item \textbf{Local} (\code{.claude/settings.local.json}) --- personal overrides, auto-gitignored
  \item \textbf{Project} (\code{.claude/settings.json}) --- team standards, committed to git
  \item \textbf{User} (\code{\textasciitilde/.claude/settings.json}) --- global personal defaults
\end{enumerate}

\subsection{Key Settings}

\begin{description}
  \item[\code{permissions}] Allow, deny, and ask rules for tool access.
    Use \code{allow} for safe operations, \code{deny} for sensitive files,
    \code{ask} for destructive actions.
  \item[\code{model}] Default model selection (opus, sonnet, haiku).
  \item[\code{hooks}] Lifecycle event handlers (see \cref{sec:hooks}).
  \item[\code{outputStyle}] Response format customization.
  \item[\code{env}] Environment variables passed to Claude sessions.
\end{description}

\subsection{Convention: Team vs.\ Personal Settings}

\practitioner{Commit \code{settings.json} for team-wide standards.
Keep \code{settings.local.json} for personal preferences.}

\begin{minted}{json}
// .claude/settings.json (committed -- team standards)
{
  "permissions": {
    "deny": ["Read(.env*)", "Read(secrets/**)"],
    "ask": ["Bash(git push:*)"]
  }
}
\end{minted}

\begin{minted}{json}
// .claude/settings.local.json (gitignored -- personal)
{
  "model": "opus",
  "outputStyle": "concise"
}
\end{minted}

\subsection{Practical Deny Rules}

\margintip{Deny rules are absolute. Unlike CLAUDE.md instructions, they cannot be overridden by conversation.}

\begin{minted}{json}
{
  "permissions": {
    "deny": [
      "Read(.env*)",
      "Read(secrets/**)",
      "Read(**/*.pem)",
      "Read(**/*.key)",
      "Read(**/credentials*)",
      "Read(**/.aws/**)",
      "Bash(rm -rf /)",
      "Bash(sudo:*)"
    ]
  }
}
\end{minted}

% -------------------------------------------------------------------
\section{Output Styles}
\label{sec:output-styles}

\marginnote[Official]{Configure via \code{outputStyle} in settings or \code{--output-style} flag.}

\official{Output styles customize how Claude structures its responses:
tone, verbosity, formatting conventions, and problem-solving protocols.}

A custom output style can encode:
\begin{itemize}
  \item Problem-solving workflows (e.g., ``always diagnose before fixing'')
  \item Parallelism strategies for tool usage
  \item Domain-specific communication patterns
  \item Formatting preferences (bullet points vs.\ prose, emoji policy)
\end{itemize}

\practitioner{Custom styles are most valuable when they encode \emph{decision protocols},
not just formatting preferences.}
A style that says ``use bullets'' saves minimal effort.
A style that says ``classify the issue type before searching for solutions''
changes the quality of every interaction.

% -------------------------------------------------------------------
\section{Visual: The Configuration Hierarchy}

\begin{figure}[H]
\centering
\begin{tikzpicture}[
  layer/.style={draw=PrimaryBlue, fill=PrimaryBlue!8, rounded corners=3pt,
    minimum width=10cm, minimum height=1.1cm, font=\small, align=center},
  label/.style={font=\footnotesize\itshape, text=PrimaryBlue!70!black},
  >=Stealth
]
  % Layers from bottom to top
  \node[layer] (output) at (0,0) {\textbf{Output Styles} --- Communication patterns, tone, problem-solving protocols};
  \node[layer] (settings) at (0,1.5) {\textbf{Settings} --- Permissions, model, hooks, env vars};
  \node[layer] (rules) at (0,3.0) {\textbf{Rules} --- Conditional context (paths: frontmatter)};
  \node[layer] (claude) at (0,4.5) {\textbf{CLAUDE.md} --- Project knowledge, build commands, architecture};

  % Precedence arrow
  \draw[->, thick, PrimaryBlue!60] (5.5,0) -- (5.5,4.5)
    node[midway, right, label] {Higher priority};

  % Scope labels
  \node[label, anchor=east] at (-5.3,4.5) {Global + Project};
  \node[label, anchor=east] at (-5.3,3.0) {Per-module};
  \node[label, anchor=east] at (-5.3,1.5) {4-level cascade};
  \node[label, anchor=east] at (-5.3,0) {Per-user};
\end{tikzpicture}
\caption{The four configuration layers, each handling a distinct concern.}
\label{fig:config-hierarchy}
\end{figure}

\begin{keyinsight}
The configuration hierarchy is designed so that each layer handles a different concern:
\code{CLAUDE.md} for project knowledge, rules for conditional context,
settings for permissions and behavior, and output styles for communication patterns.
Mixing these concerns (e.g., putting permission rules in CLAUDE.md)
reduces their effectiveness.
\end{keyinsight}
